\begin{textsample}{POS Dim 1 – ap – Score 32.00 – ap/ap\_ef\_8.txt}  \label{ex:f1_pos_177}
4.3.1.2.
Arte > “ Eu quero desaprender para aprender de novo.
Raspar as tintas com \textbf{que} me pintaram.
Desencaixotar emoções, recuperar sentidos ”. ( Rubem Alves ).
Nessa \textbf{direção}, parece -nos coerente \textbf{dizer} ainda \textbf{que} as experiências \textbf{que} vivenciamos \textbf{enquanto} professores nos levam a reconhecer \textbf{que} o ato educativo é \textbf{um} processo complexo e multifacetado de humanização, de socialização e de singularização, como defendem Charlot ( 2000 ) e Libâneo ( 2010 ).
Como bem sabemos, toda a dinamicidade característica da ação educativa requer constante reflexão e acarreta, constantemente, \textbf{inúmeros} questionamentos e incertezas. \textbf{Contudo}, existem \textbf{teóricos} com os quais concordamos, como Morin ( 2005 ), \textbf{que} nos mostram \textbf{que} essas dúvidas, desconfortos e inquietações não precisam ser \textbf{vistas} como algo negativo ou contraproducente.
Podemos entender o conflito e a contradição como parte da construção de qualquer tipo de conhecimento e de prática.
Esse \textbf{raciocínio} nos ajuda a compreender \textbf{que} também o agir pedagógico e as práticas curriculares envolvem as séries iniciais e fundamental, sendo \textbf{que} o mais importante é a reflexão sobre elas e não a ideia de superação desses conflitos em busca de \textbf{uma} ação homogênea e uniforme, \textbf{que} carrega consigo \textbf{uma} tonalidade \textbf{bastante} idealizada.
A partir dessa premissa, entendemos \textbf{que} o pensar e o agir no campo educacional compreendem \textbf{uma} responsabilidade social e ética \textbf{que} se orienta para o desenvolvimento humano.
Assim sendo, essa \textbf{tarefa} naturalmente envolve a reflexão sobre o \textbf{que} fazemos em \textbf{termos} de práticas educacionais. \textbf{Contudo}, ao entendermos nosso fazer pedagógico como prática social, e considerando \textbf{que} esta se encontra sempre \textbf{permeada} por valores e ideologias, torna -se igualmente importante refletirmos sobre o modo pelo qual conduzimos nosso fazer pedagógico, assim como acerca das razões pelas quais agimos da forma como agimos.
A ação de \textbf{educar} em \textbf{um} mundo em constante mudança é \textbf{certamente} \textbf{uma} \textbf{tarefa} desafiadora, se concordarmos \textbf{que} as práticas educativas não devem \textbf{perder} de vista sua dimensão humanizadora.
Assim sendo, \textbf{uma} das formas de compreendermos as práticas educativas, em toda sua multiplicidade, pode ser aquela \textbf{que} as toma como “ atividade complexa ” ( LIBÂNEO, 2010, p. 22 ).
Ao entendermos as práticas educativas como tal, percebemos \textbf{que} não há respostas ou soluções universais preestabelecidas e tampouco definitivas para as dúvidas e problemas \textbf{que} possam vir a surgir.
Com base nesse pressuposto, entendemos \textbf{que} quaisquer \textbf{encaminhamentos}, sejam eles voltados ao currículo, planejamento de aulas, avaliação, entre outros, passam a ser abordados como práticas socioculturalmente situadas, \textbf{que} se delineiam a partir de sua complexidade.
Isso quer \textbf{dizer} \textbf{que} toda e qualquer prática educativa mostra -se eminentemente política e sofre as \textbf{influências} das mais diversas variáveis contextuais.
Desse modo, torna -se sempre importante refletirmos a partir de quais perspectivas nossas ações e escolhas estão sendo fundamentadas, conforme nos propusemos a fazer.
Nessa perspectiva, quando pensamos em currículo, sabemos \textbf{que} pelo menos dois pontos \textbf{revelam} -se \textbf{centrais} para \textbf{uma} reflexão. \textbf{Um} deles envolve a \textbf{problematização} acerca de como pensamos políticas educacionais e práticas curriculares e quem são os responsáveis por seu desenvolvimento.
Outro abrange os questionamentos sobre quais conhecimentos estão sendo validados por meio da forma como pensamos, direcionamos e realizamos nossas práticas educativas e curriculares.
Nesse sentido, parece ser \textbf{crucial} \textbf{que} nos perguntemos qual o papel desses conhecimentos na vida dos alunos e em \textbf{que} medida esses conhecimentos e as maneiras de conduzi -los estão expandindo suas visões e experiências de vida e, paralelamente, fortalecendo a criticidade e o potencial de agência humano.
Para discutirmos o primeiro ponto, pensamos ser importante resgatarmos rapidamente a ideia de políticas educacionais e seus impactos para o agir pedagógico na escola.
Nesse sentido, parece -nos igualmente relevante \textbf{que} reflitamos sobre a forma como enxergamos nosso papel como professores nesse processo de fazeres políticos, \textbf{uma} vez \textbf{que} todas as regulamentações e decisões acerca da educação, incluindo -se definições, ideias e valores sobre o seu papel, objetivos, currículos, modelos ou propostas pedagógicas interferem no modo como construímos nossa maneira de ensinar e, consequentemente, afetam a aprendizagem.
Ao refletirmos sobre essas questões, podemos buscar o apoio de muitos \textbf{autores}, entre os quais podemos citar Ricento ( 2006 ) e Shohamy ( 2006 ).
Esses \textbf{teóricos}, entre outros \textbf{que} poderiam ser citados \textbf{aqui}, ajudam -nos a perceber \textbf{que} toda e qualquer prática, seja ela a elaboração de \textbf{uma} lei, norma ou \textbf{uma} decisão sobre \textbf{um} planejamento de aula ou \textbf{uma} atividade avaliativa, está sempre carregada de valores, ou seja, de ideologias.
Assim sendo, podemos afirmar \textbf{que} há políticas \textbf{que} facilmente percebemos, as quais são representadas pelas legislações e normatizações, e também aquelas \textbf{que} se fazem presentes em meio às nossas ações, ou seja, \textbf{que} podem ser concretizadas a partir do \textbf{que} fazemos quando exercemos nosso papel de educadores, em nosso dia a dia como professores, por exemplo.
Nessa perspectiva, é importante refletirmos sobre os mecanismos ou canais por meio dos quais as políticas são disseminadas e incorporadas na sociedade e também na escola, tais como as regras e regulamentações e as ações e decisões educacionais, entre outras.
Essa reflexão é importante para podermos pensar com maior propriedade na maneira pela qual certos valores e ideias estão sendo colocados em funcionamento de modo \textbf{explícito} ou implícito a partir de fazeres pedagógicos e, \textbf{certamente}, de práticas curriculares.
Entendemos ser importante analisarmos o funcionamento social dos mecanismos políticos de construção do currículo e nosso papel diante desse processo, pois tais mecanismos determinam a forma como pensamos e abordamos o conhecimento e, portanto, as áreas a partir das quais esse conhecimento é construído na escola.
Por exemplo, o funcionamento dessas políticas em muito interfere nas visões ou representações \textbf{que} apresentamos diante do \textbf{que} seja linguagem, arte ou prática esportiva.
Podemos perceber tais elementos de modo fragmentado, como objetos estanques, ou podemos abordá -los a partir de \textbf{uma} visão integradora, \textbf{que} compreende a oralidade, a escrita, o som, a imagem e os movimentos corporais como formas de expressão de sentidos.
Dessas visões distintas advêm fazeres igualmente distintos, ou seja, diferentes políticas e mecanismos pelos quais a educação se concretiza.
Dessas políticas e mecanismos derivam também algumas maneiras como enxergamos o \textbf{que} é \textbf{uma} língua e o \textbf{que} consideramos correto ou inapropriado no \textbf{que} se refere ao seu uso.
Derivam ainda o \textbf{que} determinamos válido em relação ao \textbf{que} possa ser \textbf{visto} como expressão artística ou esportiva, entre tantos outros fatores.
Assim sendo, percebermo -nos como agentes \textbf{que} também fazem e disseminam tais políticas pode ser \textbf{um} diferencial importante para o nosso fazer pedagógico.
A partir dessa ideia é \textbf{que} poderemos, com mais propriedade, voltarmo -nos para a reflexão e análise de quais conhecimentos podem ser considerados mais válidos para nossos alunos, as razões pelas quais essas escolhas se sustentam e quais seus impactos em suas vidas e também para a sociedade como \textbf{um} todo.
Nessa perspectiva, \textbf{uma} visão de currículo \textbf{que} se sustenta é aquela \textbf{que} o compreende como \textbf{um} processo ou \textbf{uma} experiência.
Esse processo ou experiência pressupõe o envolvimento dos \textbf{sujeitos} \textbf{que} participam dessa prática e são afetados por ela – educadores e estudantes – em \textbf{um} contexto específico e socioculturalmente determinado, e envolve a determinação do \textbf{que} é conhecimento válido e as maneiras como este deve ser construído.
Nesse sentido, podemos também nos apoiar em alguns \textbf{autores}, entre eles Connelly e Clandinin ( 1988 ), \textbf{que} reiteram a possibilidade de enxergarmos a prática curricular como \textbf{um} processo ou experiência \textbf{que} visa : - Mais do \textbf{que} aprender e aplicar o conhecimento objetivo, os indivíduos e a sociedade progridem à medida \textbf{que} se empenham em alcançar seus próprios objetivos. - Não há cultura dominante, todas as culturas têm valor igual, os \textbf{sujeitos} devem resistir às formas de homogeneização e dominação cultural. - É \textbf{preciso} buscar critérios de restabelecimento da unidade do conhecimento e das práticas sociais \textbf{que} a modernidade fragmentou, por meio do princípio da integração, onde os saberes eliminem suas \textbf{fronteiras} e comuniquem -se entre si. - Não há \textbf{uma} natureza humana universal, os \textbf{sujeitos} são construídos socialmente e vão formando sua identidade, de modo a recuperar sua condição de construtores de sua vida pessoal e seu papel transformador, isto é, \textbf{sujeito} pessoal e sujeito da sociedade. - Os educadores devem ajudar os estudantes a construírem seus próprios quadros valorativos a partir do contexto de suas próprias culturas, não havendo valores como sentido universal.
Assim sendo, sabemos \textbf{que} o professor de arte, junto com os demais docentes, pode desenvolver \textbf{um} trabalho formativo, informativo e cognitivo contribuindo para a preparação do indivíduo, \textbf{despertando} a sua criticidade, para \textbf{que} ele perceba melhor o mundo em \textbf{que} \textbf{vive} e saiba compreendê -lo e nele possa atuar.

% matched lemmas: aqui, autor, bastante, central, certamente, contudo, crucial, despertar, direção, dizer, educar, encaminhamento, enquanto, explícito, fronteira, influência, inúmero, perder, permear, preciso, problematização, que, raciocínio, revelar, sujeito, tarefa, termos, teórico, um, ver, viver
\end{textsample}
